\chapter{Build system: CMake tour}
\label{ch:cmake}

\begin{aside}
\maya{} is header-mostly. Linking it is one line of CMake; the
rest of the file is about your app.
\end{aside}

\section{Minimal CMakeLists}

\begin{lstlisting}
cmake_minimum_required(VERSION 3.28)
project(myapp CXX)
set(CMAKE_CXX_STANDARD 26)
add_subdirectory(maya)
add_executable(myapp main.cpp)
target_link_libraries(myapp PRIVATE maya)
\end{lstlisting}

That's the whole file. \maya{} is a CMake target you
\code{add\_subdirectory} into.

\section{Compile flags}

\maya{} requires:

\begin{itemize}[leftmargin=*]
  \item C++26 (\code{-std=c++2c} or \code{-std=c++26}).
  \item Concepts support.
  \item \code{constexpr} string literals.
\end{itemize}

GCC 14+ or Clang 18+ supplies these.

\section{Build types}

CMake offers Debug, Release, RelWithDebInfo, MinSizeRel.
For development, RelWithDebInfo gives speed plus debug
symbols.

\section{Sanitisers}

Wire ASan and UBSan via flags:

\begin{lstlisting}
target_compile_options(myapp PRIVATE
  -fsanitize=address,undefined)
target_link_options(myapp PRIVATE
  -fsanitize=address,undefined)
\end{lstlisting}

Catches memory bugs and undefined behaviour at runtime.

\section{Ninja}

\begin{lstlisting}
cmake -B build -G Ninja
ninja -C build
\end{lstlisting}

Ninja is faster than Make for incremental builds.

\section{Cross-compilation}

A toolchain file specifies the cross compiler. Example for
ARM Linux:

\begin{lstlisting}
set(CMAKE_C_COMPILER aarch64-linux-gnu-gcc)
set(CMAKE_CXX_COMPILER aarch64-linux-gnu-g++)
\end{lstlisting}

\section{Distribution}

For a binary release, link statically (\code{-static}) and
strip symbols (\code{strip myapp}). The result runs on any
glibc-compatible system of the same architecture.

\section{Packaging}

CPack generates DEB, RPM, ZIP, etc. from your CMake
project. One \code{cpack -G DEB} produces a deb file.

\section{Recap}

\begin{recap}
\begin{itemize}[leftmargin=*]
  \item \maya{} is one \code{add\_subdirectory} away.
  \item C++26-capable compiler required.
  \item Sanitisers via flag flip.
  \item Ninja for fast incremental builds.
  \item CPack for distribution packages.
\end{itemize}
\end{recap}

\section*{Workshop}

\begin{enumerate}[leftmargin=*]
  \item Add ASan to your build; fix any warnings.
  \item Switch from Make to Ninja; observe build time.
  \item Generate a deb package via CPack.
\end{enumerate}
